% Pursuing Stacks (cotes 134-2 to 134-8): archive cote 134-4, site page 70 (PDF page 71), typescript page 158.
% typescript. Draft compared: ps3.tex:3612-3666.
% AI-transcribed by Claude Opus 5.5 (claude-opus-5-5, adaptive thinking, effort high), Message Batches API, batch ps2 (msgbatch_01NVwr1Yuo8qJf5rykeGBXPK), 2026-11-12; not yet checked by eye.
\scanpage{134-4/70}
\mspage{158}

$\underline{W} = \unsure{Fl|F/}((Cat))$ -- \uline{all} arrows in (Cat) ! This circumstance seems tied
closely to the fact that, contrarily to quite unreasonable expectations,
\uline{we \unsure{definitely} do \uuline{not} get an axiomatic caracterization of weak equivalen}\xout{\unsure{c}}%
\uline{ces} , in terms of the type of \xout{conditions upon $\underline{W}$} properties of $\underline{W}$ we have
been working with so far. As soon as one stops for considering the matter
without prejudice, this appears as rather obvious. As a matter of fact,
using still cohomological invariants of topoi and categories, there are
lots of variants of the cohomological definition of ``weak equivalence'',
which will share all formal properties \over{\illegible{1}}{of} the latter we have been using
so far, and presumably a few more we hav'nt met yet. For instance,
starting with any \xout{commutative group} ring \inserted{k} (interesting cases would be $\underline{Z}$,
$\underline{Z}/n\underline{Z}$ , $\underline{Q}$ ), we may demand on a morphism of topoi
\[ f: X \hand{\longrightarrow} Y \]
to induce an \unsure{isomorphism\xout{e}} for cohomology with coefficients in k , or
with coefficients in any k-module, or with any twisted coefficients which
are k-modules -- already three candidates for a $\underline{W}$ , depending on a
given k ! We may vary still more, by taking, instead of just one
k , a whole bunch $(k_{i})$ of such, or a bunch of \hand{(}constant \deleted{grou}\unsure{\hand{)}} commutative
groups -- we \xout{think} are thinking of choices such as all rings $\underline{Z}/n\underline{Z}$ , with
possibly n being subjected to be prime to a given set of primes, along
the lines of the Artin-Mazur theory of ``localization'' of homotopy types.
And we may combine this with an isomorphism requirement on twisted non
commutative 1-cohomology, as for the usual notion of weak equivalence.
Also, in all the isomorphism requirements, we may restrict to cohomolo-%
gy up to a certain dimension (which will give rise to ``truncated homotopy
types\hand{''}). The impression that goes with the evocation of all these exam-%
ples, is that the theory we have been pursuing, to come to an understan-%
ding of ``models for homotopy types'', while we started with just usual
homotopy types in mind and a corresponding tacit prejudice, is a great
deal richer than what we had in mind. Yesterday's (or rather last
night's) embarassment of finding out finally I had been very silly, is
a typical illustration of the embarassment we feel, whenever a forebo-%
ding appears of our sticking to unadequate ideas\over{,}{;} still more so if it is
not just mathematics but ideas about ourselves say or about something
in which we are strongly personally implied. This embarassment then
comes as a rescue, to bar the way to an unwished for overwhelming
richess dormant in ourselves, \xout{and} ready to wash away forever those ideas
so dear to us ...

\margin{70}
